\providecommand{\kBKM}{\kappa_{\mathrm{BKM}}}
\providecommand{\kcat}{\kappa_{\mathrm{cat}}}
\providecommand{\kch}{\kappa_{\mathrm{ch}}}
\providecommand{\kchHeis}{\kappa_{\mathrm{ch}}^{\mathrm{Heis}}}
\providecommand{\kfiber}{\kappa_{\mathrm{fiber}}}
\providecommand{\PhiCHL}[1]{\Phi_{#1}^{\mathrm{CHL}}}

\section*{EOT Table 1 and Borcherds Weight Data}

\paragraph{EOT Table 1.}
Eguchi, Ooguri, and Tachikawa identify Table 1 in Appendix A of
\texttt{arXiv:1004.0956} as the character table of the Mathieu group
\(M_{24}\). Its row labels are the irreducible representation dimensions
\(1,23,45,\ldots,10395\), and its column labels are the \(26\)
conjugacy classes \(1A,2A,\ldots,23B,12A\). The coefficient table in
the same paper is equation (1.12), where the \(A_n\) record the
\(\mathcal N=4\) decomposition of the K3 elliptic genus.

A claimed extraction of \(c_N(0)\) from EOT Table 1 is a category error:
the table entries are character values. Mathieu character data become
Jacobi-form data only after a conjugacy class, and for simultaneous
twining a commuting pair, a multiplier, and a root datum have been fixed.

\paragraph{Three constants.}
The letter \(c\) is attached to three distinct constructions.
\[
\begin{array}{c|c|c}
\text{quantity} & \text{definition} & \text{role}\\ \hline
\chi_{24}(g) & \operatorname{Tr}_{\mathbf 1\oplus\mathbf {23}}(g)
& K3 twined Euler number\\
c_N^{\mathrm{CHL}}(0) & [\zeta^0q^0]\phi^{\mathrm{CHL}}_{0,1;N}
& \text{primitive CHL Borcherds input}\\
c_{\mathrm{dd}}(0;t,N) & 2\,\operatorname{wt}(F_{t,N})
& \text{diagonal-divisor catalogue constant}
\end{array}
\]
The first line is representation-theoretic. The second line fixes the
primitive CHL denominator. The third line belongs to the Gritsenko and
Clery diagonal-divisor catalogue.

\paragraph{Primitive CHL denominator.}
For \(N\in\{1,2,3,4,6\}\), the CHL denominator row is
\[
\begin{array}{c|c|c}
N & c_N^{\mathrm{CHL}}(0) & \kBKM(\PhiCHL{N})\\ \hline
1 & 10 & 5\\
2 & 8 & 4\\
3 & 6 & 3\\
4 & 4 & 2\\
6 & 2 & 1
\end{array}
\]
Borcherds' weight theorem gives, row by row,
\[
  \kBKM(\PhiCHL{N})
  =
  \operatorname{wt}(\PhiCHL{N})
  =
  \frac{c_N^{\mathrm{CHL}}(0)}2 .
\]
At \(N=1\), the Eichler and Zagier generator has
\[
  \phi_{0,1}\big|_{q^0}
  =
  \zeta^{-1}+10+\zeta,
\]
hence \(c_1^{\mathrm{CHL}}(0)=10\) and
\(\kBKM(\Delta_5)=5\).

\paragraph{The \(K3\times E\) invariant split.}
The compact \(K3\times E\) datum carries five named readings:
\[
\begin{array}{c|c|l}
\text{invariant} & \text{value} & \text{construction}\\ \hline
\kcat(K3\times E) & 0
& \chi(\mathcal O_{K3})\chi(\mathcal O_E)=2\cdot0\\
\kch(K3\times E) & 0
& \text{compact Hodge supertrace}\\
\kchHeis(K3\times E) & 3
& \text{Heisenberg Mukai chiral specialisation}\\
\kBKM(\mathfrak g_{\Delta_5}) & 5
& \text{Borcherds denominator weight of }\Delta_5\\
\kfiber(K3) & 24
& \text{Mukai-lattice rank of the K3 fibre}
\end{array}
\]
The Borcherds entry is the Hall-Drinfeld/Borcherds branch. The K3
self-mirror Yangian branch uses a separate Mukai-lattice input.

\paragraph{Diagonal-divisor catalogue.}
Gritsenko and Clery classify the dd-modular forms by triples
\((t,N;k)\). The row data are
\[
\begin{array}{c|c|c}
(t,N;k) & \operatorname{wt}(F_{t,N}) & c_{\mathrm{dd}}(0;t,N)\\ \hline
(1,1;5) & 5 & 10\\
(2,1;2) & 2 & 4\\
(1,2;3) & 3 & 6\\
(3,1;1) & 1 & 2\\
(1,3;2) & 2 & 4\\
(4,1;\tfrac12) & \tfrac12 & 1\\
(1,4;\tfrac32) & \tfrac32 & 3\\
(2,2;1) & 1 & 2
\end{array}
\]
The common row with the CHL table is \((t,N;k)=(1,1;5)\), where
\(F_{1,1}=\Delta_5\). The remaining rows have their own paramodular
groups, characters or multiplier systems, and Borcherds products.

\paragraph{Excluded row.}
The five-entry row
\[
\begin{array}{c|ccccc}
N & 1 & 2 & 3 & 4 & 6\\ \hline
c(0) & 10 & 4 & 2 & 2 & 2
\end{array}
\]
requires its own Jacobi input and product normalisation. Substituting it
into the primitive CHL denominator theorem would give weights
\[
\begin{array}{c|ccccc}
N & 1 & 2 & 3 & 4 & 6\\ \hline
c(0)/2 & 5 & 2 & 1 & 1 & 1
\end{array}
\]
which is a different Borcherds row from the CHL table above.

\paragraph{Conclusion.}
EOT Table 1 supplies Mathieu character values. The primitive CHL
Borcherds constants are the five \(c_N^{\mathrm{CHL}}(0)\) displayed
above, and the Gritsenko and Clery constants are the eight
triple-indexed \(c_{\mathrm{dd}}(0;t,N)\). A valid comparison must name
the table, Jacobi input, cover or multiplier, and product normalisation.
